\subsection{Rappresentazione di Heisenberg}
\[
A^H(t) = S^{-1}(t, t_0) A^S(t_0) S(t, t_0), \quad A^S(t_0) \text{ costante nel tempo.}
\]
Per \( \ket{\psi} \) ho \( \ket{\psi(t)} = S(t, t_0) \ket{\psi(t_0)} \) e ho elemento di matrice
\begin{align*}
\braket{\psi_H(t_0)|A^H(t)|\psi_H(t_0)} &= \braket{\psi_S(t_0)|S^{-1}(t, t_0) A^S(t_0) S(t, t_0)|\psi_S(t_0)} = \\
&= \braket{\psi_S(t)|A^S(t_0)|\psi_S(t)}
\end{align*}
Quindi passo da \( S \) ad \( H \) e nel primo l'operatore \`e costante nel tempo ed evolvono gli stati, nel secondo \`e il contrario.

Allora posso prendere gli \( \ket{n(t)} \) (fissi nel tempo) e uso
\[
A^H(t) \ket{n(t)} = \lambda_n \ket{n(t)}
\]
\[
S^{-1}(t, t_0) A^S(t_0) S(t, t_0) \ket{n(t)} = \lambda_n \ket{n(t)}
\]
\[
A^S(t_0) \underbrace{S(t, t_0) \ket{n(t)}}_{\ket{n(t_0)}} = \lambda_n \underbrace{S(t, t_0) \ket{n(t)}}_{\ket{n(t_0)}}
\]
\[
\ket{n_H(t)} = S^{-1}(t, t_0) \ket{n_S(t_0)}
\]
Quindi una generica misura
\[
\braket{n_S(t_0)|\psi_S(t)} = \braket{n_S(t_0)|S(t, t_0)|\psi_S(t_0)} = \braket{n_H(t)|\psi_H(t)}.
\]

Per studiare l'evoluzione temporale di \( A^H(t) \) faccio
\begin{align*}
\frac{d}{dt} A^H(t) &= \frac{d}{dt} \left( S^{-1}(t, t_0) A^S(t_0) S(t, t_0) \right) = \\
&= \frac{\partial S^{-1}}{\partial t} A^S(t_0) S + S^{-1} A^S \frac{\partial S}{\partial t} + S^{-1} \frac{\partial A^S}{\partial t} S,
\end{align*}
ma so che
\[
i \hbar \frac{\partial S}{\partial t} = \mathcal{H} S \; \Rightarrow \; -i \hbar \frac{\partial S^\dag}{\partial t} = S^\dag \mathcal{H}^\dag \; \Rightarrow \; -i \hbar \frac{\partial S^{-1}}{\partial t} = S^{-1} \mathcal{H}
\]
quindi 
\[
\frac{d}{dt} A^H(t) = -\frac{1}{i \hbar} S^{-1} \mathcal{H} A^S S + S^{-1} A^S \frac{\mathcal{H} S}{i \hbar} + \frac{\partial}{\partial t} A^H(t).
\]
Poich\'e per definizione \( S(t, t_0) A^H(t) = A^S(t_0) S(t, t_0) \), ho
\begin{align*}
\frac{d}{dt} A^H(t) &= -\frac{1}{i \hbar} S^{-1} \mathcal{H} S A^H + \frac{1}{i \hbar} A^H S^{-1} \mathcal{H} S + \frac{\partial}{\partial t} A^H(t) = \\
&= \frac{1}{i \hbar} \left[ A^H, \underbrace{S^{-1} \mathcal{H} S}_{\mathcal{H}^H} \right] + \frac{\partial}{\partial t} A^H(t)
\end{align*}
\begin{align*}
\frac{d}{dt} A_{ij}(t) &= \frac{d}{dt} \braket{\psi_i^S(t)|A^S(t_0)|\psi_j^S(t)} = \\
&= \frac{d}{dt} \braket{\psi_i^H(t_0)|A^H(t)|\psi_j^H(t_0)} = \\
&= \frac{1}{i \hbar} \braket{\psi_i^H(t_0)|\left\{ \left[ A^H(t), \mathcal{H}^H \right] + \frac{\partial}{\partial t} A^H(t) \right\}|\psi_j^H(t_0)}
\end{align*}
per\`o \( A^H = S^{-1} A^S S \) e \( A^H \mathcal{H}^H = S A^S S S^{-1} \mathcal{H} S = S A^S \mathcal{H} S \)
\[
\frac{d}{dt} A_{ij}(t) = \braket{\psi_i^S(t)|\frac{1}{i \hbar} \left[ A^S(t_0), \mathcal{H}^S \right]|\psi_j^S(t)} + \braket{\psi_i^S(t)|\frac{\partial A^S}{\partial t}|\psi_j^S(t)}.
\]

Esempio: \( A \) e \( B \) costanti del moto \( \Rightarrow [A, B] \) costante del moto. Infatti
\[
[A, \mathcal{H}] = [B, \mathcal{H}] = 0 \quad \Rightarrow \quad \left[ [A, B], \hat{\mathcal{H}} \right] = 0.
\]